\begingroup
\section{Foundation F148: the exact cumulative-free material matrix}
\label{module:F148}
\noindent\textit{Audited source: \nolinkurl{convex_largeN_polya_szego_stage148_NR_quadrature_primitive_material_matrix.tex}.}
\subsection{Clausen material derivative as a length gradient}

Let $I_j=[\theta_j,\theta_{j+1}]$ be a cyclic interval partition of
$\T$, with length $h_j$ and midpoint $c_j$.  Put
\begin{equation}\label{F148:eq:clausen}
 \Phi(x)=\sum_{n\ge1}\frac{\sin(nx)}{n^2},\qquad
 K(x)=\Phi'(x)=-\log\left|2\sin\frac x2\right|.
\end{equation}
The sampled conjugate of the quadratic fan spline is
\begin{equation}\label{F148:eq:w}
 w_i=\frac1\pi\sum_jh_j\Phi(c_i-c_j).
\end{equation}
Along a positive affine path, let $h_j'=d_j$, $\sum_jd_j=0$, and let
$s_j=c_j'$.  Stage 144 gives
\begin{equation}\label{F148:eq:material-old}
 W_i:=\dot w_i=\frac1\pi\sum_{j\ne i}
 \{d_j\Phi(c_i-c_j)+h_j(s_i-s_j)K(c_i-c_j)\}.
\end{equation}

Fix the target $i$.  Relabel cyclically and choose a partition endpoint
$\theta_0$ as a cut; write $\theta_0+2\pi$ for its second copy.  The cut
may be chosen at circular distance at least $\pi/3$ from $c_i$.  This
target-dependent relabelling changes neither \eqref{F148:eq:w} nor
\eqref{F148:eq:material-old}.

For $k<i$ define
\begin{align}
 \mathscr E_{ik}^-&=
 \sum_{j<k}h_jK(c_i-c_j)+\frac{h_k}{2}K(c_i-c_k)
 -\int_{\theta_0}^{c_k}K(c_i-y)\,dy. \label{F148:eq:Eminus}
\end{align}
For $k>i$ define
\begin{align}
 \mathscr E_{ik}^+&=
 \int_{c_k}^{\theta_0+2\pi}K(c_i-y)\,dy
 -\frac{h_k}{2}K(c_i-c_k)-\sum_{j>k}h_jK(c_i-c_j).
 \label{F148:eq:Eplus}
\end{align}
The two arcs in these definitions do not contain $c_i$ in their
interiors.  Finally put
\begin{align}
 E_i^-&=\sum_{j<i}h_jK(c_i-c_j)
       -\int_{\theta_0}^{\theta_i}K(c_i-y)\,dy,\notag\\
 E_i^+&=\sum_{j>i}h_jK(c_i-c_j)
       -\int_{\theta_{i+1}}^{\theta_0+2\pi}K(c_i-y)\,dy,
 \label{F148:eq:E-self}
\end{align}
and define the matrix
\begin{equation}\label{F148:eq:matrix}
 \mathscr E_{ik}=\begin{cases}
 \mathscr E_{ik}^-,&k<i,\\
 (E_i^--E_i^+)/2,&k=i,\\
 \mathscr E_{ik}^+,&k>i.
 \end{cases}
\end{equation}

\begin{theorem}[Exact cumulative-free material matrix]
\label{F148:thm:matrix}
For every zero-sum length velocity $d$,
\begin{equation}\label{F148:eq:W-matrix}
 \boxed{W_i=\frac1\pi\sum_k d_k\mathscr E_{ik}.}
\end{equation}
Thus no node velocity $s_j$, root velocity, or cumulative sum of the
$d_j$ occurs in the final coefficient matrix.
\end{theorem}

\begin{proof}
Temporarily keep the cut fixed.  Its length coordinates give
\begin{equation*}
 \frac{\partial c_j}{\partial h_k}=
 \begin{cases}0,&j<k,\\1/2,&j=k,\\1,&j>k.
 \end{cases}
\end{equation*}
If $M_{ik}=\pi\,\partial w_i/\partial h_k$, differentiation of
\eqref{F148:eq:w} gives
\begin{equation}\label{F148:eq:M}
 M_{ik}=\Phi(c_i-c_k)+\sum_jh_j
 \left(\frac{\partial c_i}{\partial h_k}
       -\frac{\partial c_j}{\partial h_k}\right)K(c_i-c_j),
\end{equation}
where the self coefficient of $K(0)$ is zero before the Abel limit.

For $k<i$, equation \eqref{F148:eq:M} becomes
\begin{equation*}
 M_{ik}=\Phi(c_i-c_k)+\sum_{j<k}h_jK(c_i-c_j)
                    +\frac{h_k}{2}K(c_i-c_k).
\end{equation*}
The fundamental theorem of calculus gives exactly
\begin{equation*}
 \int_{\theta_0}^{c_k}K(c_i-y)\,dy
 =\Phi(c_i-\theta_0)-\Phi(c_i-c_k).
\end{equation*}
After rotating the origin to the cut, the first term on the right is
$\Phi(c_i)$.  Hence
$M_{ik}=\Phi(c_i)+\mathscr E_{ik}^-$.  The same calculation on the
right arc gives
$M_{ik}=\Phi(c_i)+\mathscr E_{ik}^+$ for $k>i$.

For $k=i$, \eqref{F148:eq:M} reads
\begin{equation*}
 M_{ii}=\frac12\sum_{j<i}h_jK(c_i-c_j)
       -\frac12\sum_{j>i}h_jK(c_i-c_j).
\end{equation*}
The exact integrals over the two sides of $I_i$ give
$M_{ii}=\Phi(c_i)+(E_i^--E_i^+)/2$.  Therefore
\begin{equation*}
 M_{ik}=\Phi(c_i)+\mathscr E_{ik}\qquad\hbox{for every }k.
\end{equation*}
Since $\sum_kd_k=0$, the common Clausen term disappears and
\eqref{F148:eq:W-matrix} follows.  This also proves that the result is
independent of the chosen cut.
\end{proof}

\subsection{The matrix is a cotangent integral of a cellwise sawtooth}

On each cell define
\begin{equation}\label{F148:eq:eta}
 \eta(y)=\begin{cases}
 -(y-\theta_j),&\theta_j<y<c_j,\\
 \theta_{j+1}-y,&c_j<y<\theta_{j+1}.
 \end{cases}
\end{equation}
Then, distributionally on $I_j$,
\begin{equation}\label{F148:eq:eta-properties}
 d\eta=h_j\delta_{c_j}-\mathbf1_{I_j}\,dy,\qquad
 |\eta|\le h_j/2,\qquad
 \int_{I_j}\eta=0.
\end{equation}
Moreover
\begin{equation}\label{F148:eq:Kprime}
 K'(x)=-\frac12\cot\frac x2=-\pi\kappa(x),
 \qquad \kappa(x)=\frac1{2\pi}\cot\frac x2.
\end{equation}

\begin{proposition}[Exact sawtooth representation]
\label{F148:prop:sawtooth}
For $k<i$ and $k>i$, respectively,
\begin{align}
 \mathscr E_{ik}^-&=
 \int_{\theta_0}^{c_k}\eta(y)K'(c_i-y)\,dy,
 \label{F148:eq:saw-left}\\
 \mathscr E_{ik}^+&=-
 \int_{c_k}^{\theta_0+2\pi}\eta(y)K'(c_i-y)\,dy.
 \label{F148:eq:saw-right}
\end{align}
The terminal half-atom at $c_k$ closes the cumulative mass to zero.
Likewise $E_i^-$ and $E_i^+$ are the corresponding full-cell integrals
on the two sides of $I_i$.  Hence every matrix entry is assembled from
the same cellwise mean-zero sawtooth and one cotangent frequency.
\end{proposition}

\begin{proof}
On the left arc, the signed midpoint-quadrature measure consists of all
full atoms before $I_k$, half of the atom at $c_k$, and minus Lebesgue
measure on the arc.  Its cumulative distribution is exactly $\eta$ and
vanishes at both endpoints.  Stieltjes integration by parts therefore
gives
\begin{equation*}
 \int K(c_i-y)\,d\eta(y)
 =-\int\eta(y)\frac d{dy}K(c_i-y)\,dy
 =\int\eta(y)K'(c_i-y)\,dy,
\end{equation*}
which is \eqref{F148:eq:saw-left}.  The right arc has the opposite orientation,
giving \eqref{F148:eq:saw-right}.  The full-side statements are identical.
\end{proof}

\begin{remark}[Both former cancellations are now identities]
The endpoint vanishing in stage 147 is encoded by the fact that $\eta$
is zero at every cell endpoint.  The cumulative-velocity cancellation is
encoded by $\int_{I_j}\eta=0$ and the disappearance of the common
$\Phi(c_i)$ in Theorem~\ref{F148:thm:matrix}.  Thus neither cancellation needs
to be recovered after a triangle inequality.
\end{remark}
\endgroup
